\section{Related work}
\label{sec:related}

\paragraph{Prediction-powered inference.}
PPI \citep{angelopoulos2023ppi} combines a small labelled sample with predictions on a large unlabelled one to
obtain valid intervals for means, quantiles and regression coefficients, whatever the predictor. PPI++
\citep{angelopoulos2023ppiplus} tunes the weight placed on the predictions so that the interval is never wider
(asymptotically) than the labelled-only interval, and gives the efficiency bound and equality condition that
frame our question. Variants change what is fixed: cross-prediction-powered inference
\citep{zrnic2023cross} handles predictors trained on the same labelled data by cross-fitting; active statistical
inference \citep{zrnic2024active} chooses which records to label, and stratified PPI \citep{fisch2024stratified}
reweights strata. All of these either add information (a better predictor, a choice of labels) or change the
sampling design. We keep the labelled set, the unlabelled set and the predictor fixed and ask only how
the available information is \emph{used}.

\paragraph{Semi-supervised and doubly-robust inference.}
Semi-supervised estimation of a mean from labelled and unlabelled covariates, with an efficiency bound,
is studied by \citet{zhang2019semisup}, and for linear regression by \citet{chakrabortty2018efficient}.
The cross-fitted, orthogonal-score construction behind our adjusters is that of double/debiased machine learning
\citep{chernozhukov2018dml}. From this view PPI++ is the special case in which the nuisance function is restricted
to be affine in $f$; the question is how much of the remaining distance to the bound can be recovered from the same labelled data,
and at what finite-sample cost. We read the abstracts of the semi-supervised papers and use them for positioning only.
